\begin{abstract}
Test-time scaling offers several ways to spend additional inference computation, but transferring these methods to molecular optimization creates a comparison problem: some systems expose a candidate pool, whereas others return one molecule. We benchmark nine inference strategies on two fixed molecular instruction-tuned models and all 3,905 task--input pairs across ten multi-property tasks. Every strategy returns one final molecule under a shared evaluator, with internal computation and property-tool use specified separately. We examine three aspects of the resulting procedures: task success, generation workload, and final-molecule quality. Sampling with a public task verifier achieves macro success rates of 81.13\% and 86.18\%. Yet returning the first candidate from the same recorded calls yields only 32.36\% and 28.04\%; verifier-assisted success equals candidate-prefix coverage by construction. Graph of Thoughts achieves lower directional success but higher success under minimum-improvement thresholds. These findings distinguish successful candidate generation from usable output selection, and directional improvement from improvement magnitude. The benchmark characterizes declared molecular adaptations of prior inference mechanisms, rather than establishing an equal-compute or universal ranking of test-time scaling methods.
\end{abstract}
